Automated harness optimization improves the code that wraps a language model
into an agent through an iterative propose--evaluate loop. Recent systems such
as Meta-Harness replace compressed textual feedback with a coding agent that
can inspect the full filesystem history of prior candidates. This gives the
proposer access to raw code, traces, and diffs, but full access is not
necessarily useful access. As the optimization horizon grows, the expanding
history can dilute the proposer's attention and force it to rediscover which
past runs and files matter for the next edit. We formulate this bottleneck as
an \emph{online context curation} problem: before each proposal, the optimizer
must choose which prior iterations, trace detail, and workspace files should be
made salient under a limited inspection budget. We introduce \name, a framework
that treats this choice as a control variable in automated harness
optimization. \name combines two composable curation strategies. The
iteration-level strategy bounds and labels historical evidence, beginning with
narrow best/worst references and expanding when optimization stalls. This
strategy alone improves held-out performance in the main comparisons. The
file-level strategy can be added to assign delayed credit to files read by the
proposer and to prioritize files whose inspection has correlated with
productive edits. The two strategies are not independent additive components:
both route the same bounded proposer attention, but at different granularities.
Across three benchmarks spanning software engineering and long-horizon memory,
\name improves over the uncurated full-context baseline in most (benchmark,
proposer) settings and introduces lower-cost operating points on the
cost--quality frontier. On SWE-bench Verified, the strongest curated harness
raises full-set passrate from 0.440 to 0.640 over the source harness. These
results suggest that proposer context, usually treated as a static prompt
artifact, is itself an online optimization variable that shapes harness
search efficiency.
